% Full reconstruction of the audited structural slice/finite-descent proof.
\begin{proposition}[Proved: Reduction of a complete local puncture to surfaces]
\label{D:local-surface-reduction}
Let $(A,\mathfrak m,k)$ be a complete excellent normal local domain of
dimension $d\ge2$ with algebraically closed residue field, and let
$U=\operatorname{Spec}A\setminus\{\mathfrak m\}$. There are finitely
many complete normal local surface domains $B_1,\ldots,B_t$, with residue
field $k$, and an integer $e\ge0$ for which there is a continuous
surjection
\[
 \left(\mathop{\ast}_{i=1}^t
 \pi_1(\operatorname{Spec}B_i\setminus\{\mathfrak m_i\})\right)
 \ast\widehat F_e\ \longrightarrow\ \pi_1(U).
\]
The star denotes free profinite product. In equicharacteristic $p>0$,
all the $B_i$ are in the same equicharacteristic setting. Thus any
rank-linear $H^1$ bound, fixed-rank exponential defining-characteristic
representation bound, or uniform quadratic representation bound for
complete normal local surface punctures passes to $U$.
\end{proposition}
\begin{proof}
Choose $d-2$ elements of a system of parameters, let $I$ be their ideal,
and set $C=(A/I)_{\rm red}$. Excellence and catenarity show that every
minimal component of $C$ has dimension two. Write
$W=\operatorname{Spec}C\setminus\{\mathfrak m_C\}$.

A connected finite etale cover of $U$ normalizes to a finite complete
local normal domain $B$ over $A$. Its dimension is $d$ and its
connectedness dimension is $d$, since it is a domain. The ideal $IB$
has cohomological dimension at most $d-2$. The connectedness theorem
\citep[Tag 0ECR, with Tag 0DX9]{C_StacksLocalSlice} therefore shows
that the punctured closed subscheme cut out by $IB$ is connected:
the relevant strict inequality is $d-2<d-1$. This is precisely the
pullback of the given cover to $W$, up to nilpotents. Consequently
$\pi_1(W)\to\pi_1(U)$ is surjective.

Normalize $C$. Its normalization is finite, and it is a disjoint union
of complete normal local surface domains $B_i$; the residue fields are
$k$ because they are finite over the algebraically closed field $k$.
Their punctures give a finite surjection $W'\to W$. We now explain the
group-theoretic content of finite-proper descent. Choose basepoints on
the finitely many connected components of $W'$, and one transport for
each connected component of $W'\times_W W'$. Effective descent of finite
etale covers is \citet[Expose IX, Theorem 4.7]{C_SGA1}. The restrictions
of a cover to the vertex components, together with these finitely many
transports, determine its descent data: an isomorphism between finite
etale covers on a connected scheme is determined by its value at one
geometric point. Choosing a spanning tree in the finite component graph
leaves only finitely many independent transport generators.

More formally the resulting homomorphism from the free profinite product
of the vertex groups and a finite-rank free profinite group has dense
image: an invariant subset in a finite fibre is invariant under all
vertex monodromies and all transports, hence defines compatible
subcovers upstairs and descends to a subcover downstairs. A connected
cover admits no proper such subset. Compactness turns dense image into
surjectivity. This proves the displayed map, and also shows why mere
normal generation would not have been enough.

For completeness, on $H^1$ a cocycle is controlled by its vertex
restrictions and at most $er$ transport coordinates. For representation
classes over $\mathbf F_q$, choose a representative on every vertex;
relative conjugations and transports contribute at most
$|\operatorname{GL}_r(\mathbf F_q)|^{e+t-1}$. Multiplying the finitely
many surface bounds therefore preserves both the fixed-rank
$q^{c_r}$ form and the quadratic $q^{cr^2}$ form. No quantitative
mixed-characteristic assertion is made here.
\end{proof}
